% =====================================================================
%  Supplemental Material
% =====================================================================
\documentclass[aps,prl,onecolumn,superscriptaddress,amsmath,amssymb,floatfix]{revtex4-2}
\usepackage{graphicx,bm,xcolor,booktabs}
\usepackage[colorlinks=true,citecolor=blue!50!black,linkcolor=blue!50!black]{hyperref}
\newcommand{\gd}{\dot\gamma}
\newcommand{\ts}{\tau^{\star}}
\newcommand{\kT}{k_BT}
\newcommand{\todo}[1]{\textcolor{red}{[#1]}}
\renewcommand{\theequation}{S\arabic{equation}}
\renewcommand{\thetable}{S\arabic{table}}
\renewcommand{\thefigure}{S\arabic{figure}}

\begin{document}
\title{Supplemental Material: Unified kinetic description of flow stress, strain-rate sensitivity, dynamic strain aging, and creep in body-centered-cubic metals}
\author{Matthew Maron}
\author{Giacomo Po}
\author{Yang Li}
\maketitle

\section{S1. Components of the flow rule}

\subsection{Glide}
The mobile density moves with the kink-pair velocity
\begin{equation}
v_g=\frac{\ts b}{B_d}\exp\left\{-\frac{\Delta H_0}{2\kT}\max\left[\left(1-\left(\frac{\ts}{\tau_P}\right)^p\right)^q-\frac{T}{T_0},\,0\right]\right\}.
\end{equation}
For prismatic slip in Zr the drag coefficient interpolates smoothly between the kink drag $B_k$ and the free-flight drag $B_0$ as the barrier vanishes.

\subsection{Dislocation density}
The total density obeys
\begin{equation}
\frac{d\rho}{d\gamma}=k_1(\sqrt\rho+\Lambda_b)-k_2\rho-\frac{K D(T)}{\gd}\rho^2,
\end{equation}
with $k_2=k_{20}\exp(q_2/\kT)$ and $D=D_0\exp(-Q_D/\kT)$. The forest and mobile densities are $\rho_f=\beta\rho$ and $\rho_m=f_m\rho_f$. The flow stresses reported in the letter use the steady state, which is the root of
\begin{equation}
k_1(x+\Lambda_b)-k_2x^2-\frac{K D}{\gd}x^4=0,\qquad x=\sqrt\rho .
\end{equation}
The left-hand side is positive at $x=0$ and concave for $x>0$, so the positive root is unique. When the last term dominates, $x\propto(\gd/KD)^{1/3}$ and $\tau_a\propto(\gd/D)^{1/3}$. The boundary term is $\Lambda_b=\sum_ic_i/d_i$ over laths, blocks, prior-austenite grains and particles.

\subsection{Solute aging}
Arrested segments wait at forest junctions for $t_a=\rho_mb\lambda_f/\gd$ with $\lambda_f=\rho_f^{-1/2}$. The atmosphere concentration is
\begin{equation}
\frac{C}{C_0}=1+\sum_iC_i\left[1-e^{-x_i}\right],\qquad x_i=\left(t_a\nu_ie^{-Q_i/\kT}\right)^{\alpha_i},
\end{equation}
and the Taylor stress is multiplied by $s_c=\sqrt{C/C_0}$, so that $\tau_a=\alpha\mu b\,s_c(\sqrt{\rho_f}+\Lambda_b)$.

\subsection{Particle pinning}
Particles of areal density $\lambda_p^{-2}$ are met $n_p=\lambda_p^{-1}$ times per unit glide distance. The particle stress $\tau_p$ solves
\begin{equation}
\gd=\frac{\rho_mb}{n_p}\left(\nu_{\rm th}+\nu_{\rm cl}\right).
\end{equation}
Thermal detachment and glide-assisted local climb proceed at
\begin{align}
\nu_{\rm th}&=\nu_0\exp\left\{-\frac{F}{\kT}\left[1-\left(\frac{\tau_p}{\hat\tau}\right)^p\right]^q\right\},\\
\nu_{\rm cl}&=\frac{v_c}{h_{\rm eff}},\qquad
v_c=\frac{2\pi D_c}{b\ln(R/r_c)}\left[\exp\left(\frac{\tau_p\Omega}{\kT}\right)-1\right],\qquad
h_{\rm eff}=h\max\left(1-\frac{\tau_p}{\hat\tau},10^{-3}\right),
\end{align}
with $\hat\tau=\hat\tau_0\,\mu(T)/\mu_0$. The reduction of the climb height describes a line pressed against the particle, as in local-climb models~\cite{S_BrownHam}, and produces threshold-type creep. Both rates are non-decreasing in $\tau_p$, so $\tau_p$ is unique.

\subsection{Diffusional flow}
Nabarro--Herring and Coble flow add the rate
\begin{equation}
\gd_{\rm diff}=42\,\frac{\Omega\tau}{\kT d^2}\left(D_v+\frac{\pi\delta}{d}D_b\right).
\end{equation}

\subsection{Uniqueness and constant-stress evaluation}
With the state fixed by $(\gd,T)$, $v_g$ is strictly increasing in $\ts$ and $\tau_p$ is unique, so Eq.~(1) of the letter has exactly one solution for every imposed rate. The function $\tau(\gd)$ may decrease where aging gives $m<0$, but it is always single valued. For creep and for the maps, $\tau(\gd)$ is tabulated on $10^{-14}\le\gd\le10^{6}$~s$^{-1}$ at each temperature and inverted; where $\tau(\gd)$ is non-monotonic the lowest-rate branch that reaches the applied stress is taken, and $\gd_{\rm diff}$ is added.

\section{S2. Derivation of the rate-sensitivity decomposition}
Write $\tau=\ts+R$, where $\gd=G(\ts;\bm s)=\rho_mb\,v_g(\ts)$ is the intrinsic mobility and $R=\tau_a(\bm s)+\sum_g\tau_g$ collects all resistances. Along the solution,
\begin{equation}
d\ln\gd=A\,d\ts+\frac{\partial\ln G}{\partial\bm s}\cdot\frac{d\bm s}{d\ln\gd}\,d\ln\gd,\qquad A\equiv\frac{\partial\ln G}{\partial\ts},
\end{equation}
so that
\begin{equation}
\frac{d\tau}{d\ln\gd}=\frac1A+B,\qquad
B\equiv\frac{dR}{d\ln\gd}-\frac1A\,\frac{\partial\ln G}{\partial\bm s}\cdot\frac{d\bm s}{d\ln\gd}.
\end{equation}
With $\Lambda=AB$ this gives
\begin{equation}
m=\frac1\tau\frac{d\tau}{d\ln\gd}=\frac{1+\Lambda}{\tau A},\qquad V^\star=\frac{\kT}{m\tau}=\frac{\kT A}{1+\Lambda}.
\end{equation}
The mobility sensitivity is
\begin{equation}
A=\frac1\ts+\frac{\Delta H_0pq}{2\kT\tau_P}\,x^{p-1}(1-x^p)^{q-1},\qquad x=\frac{\ts}{\tau_P}.
\end{equation}
The contributions to $B$ are as follows. For the particle stress, $d\tau_p/d\ln\gd=1/A_p$ with $A_p=\sum_k(\nu_k/\nu)A_k$ and
\begin{align}
A_{\rm th}&=\frac{Fpq}{\kT\hat\tau}\,y^{p-1}(1-y^p)^{q-1},\qquad y=\frac{\tau_p}{\hat\tau},\\
A_{\rm cl}&=\frac{\Omega}{\kT}\frac{e^u}{e^u-1}+\frac{1}{\hat\tau-\tau_p},\qquad u=\frac{\tau_p\Omega}{\kT}.
\end{align}
Diffusion-controlled recovery lowers $\rho$ as $\gd$ decreases and therefore contributes $B>0$. Aging gives, with $t_a\propto\gd^{-1}$,
\begin{equation}
\frac{d\ln(C/C_0)}{d\ln\gd}=-\frac{\sum_iC_i\alpha_ix_ie^{-x_i}}{C/C_0}<0,
\end{equation}
so its contribution to $B$ is negative, and $m<0\iff\Lambda<-1$. In practice $A$ is evaluated analytically at the solved $\ts$, $d\tau/d\ln\gd$ is obtained by differentiating the full solution, and $B$ follows from its definition. The identity agrees with finite differences to better than $10^{-5}$ for all three materials, including points with $m<0$.

\section{S3. Material inputs and calibration}
Conventions: $\gd=\dot\epsilon/M_s$ and $\sigma=\tau/M_s$, with $M_s=0.333$ (W, Eurofer-97) and $0.45$ (Zr). The Zr and Eurofer-97 axial flow stresses are compared as $\sigma/3$. The global $m$ and $V^\star$ are least-squares fits over the experimental rate triplets.

The parameters were determined in the order of the hierarchy. The kink-pair parameters follow from the low-temperature, high-rate data; the athermal and storage parameters from the plateau; the aging parameters from the position of the stress hump and its shift with rate; and the recovery and climb parameters from the high-temperature, low-rate data. For W only level I parameters are required in the tensile range. For Zr the level II parameters were fitted to the resolved shear stresses of Derep \emph{et al.} at two rates, with a weak weight on $m$. For Eurofer-97 all parameters were fitted jointly to the flow stresses of Vanaja \emph{et al.} (three rates, 300--873~K) and to the minimum creep rates of Refs.~\cite{S_Fern,S_Yu}. The objective was the sum of squared relative stress errors plus $0.01$ times the sum of squared $\log_{10}$ rate errors, minimized by bounded Powell iterations. The activation energy for particle climb was set equal to that for forest recovery. The resulting mean flow-stress error is 4.1\%, and the creep rms error is 0.55 decades, with 52 of 54 rates within one decade. The diffusional-flow inputs were not fitted: $d=21~\mu$m (prior-austenite grain size), $D_v$ equal to the recovery diffusivity, and $D_b=2\times10^{-4}\exp(-1.8~\mathrm{eV}/\kT)$~m$^2$/s with $\delta=2b$. For W and Zr the recovery and diffusional terms used in the maps are based on self-diffusion data and are predictions.

\begin{table}[h]
\caption{Parameters. Values marked $^\dagger$ are lumped with geometry ($h$, $\ln R/r_c$) and are not bare physical constants.}
\begin{ruledtabular}
\begin{tabular}{llll}
 & W (I) & Zr (II) & Eurofer-97 (III)\\
\hline
$b$ (nm), $\mu_0$ (GPa), $T_m$ (K) & 0.2722, 161, 3695 & 0.3233, 33, 2128 & 0.2737, 80 [$\mu(T)$], 1811\\
$\Delta H_0$ (eV), $\tau_P$ (GPa), $p$, $q$ & 2.4, 0.8, 0.83, 1.69 & 3.75, 0.30, 0.86, 1.69 & 1.49, 0.68, 0.65, 1.7\\
$T_0/T_m$ & 0.9 & 0.7 & 0.9\\
$\alpha$ & 0.28 & 0.36 & 0.296\\
$k_1$ (m$^{-1}$), $k_{20}$, $q_2$ (eV) & $5.5\times10^9$, 500, $-0.011$ & $5.5\times10^9$, 500, $-0.005$ & $7.8\times10^{10}$, 6750, $-0.007$\\
$\beta$, $f_m$ & 1, 1 & 1, 1 & 0.1, 1\\
$\Lambda_b$ (m$^{-1}$) & -- & -- & $2.74\times10^6$\\
aging $C$, $Q$ (eV), $\alpha$ & -- & 0.30, 1.80, 2/3 (O) & 1.21, 0.97, 1/3 (C, N)\\
recovery $K$, $Q_D$ (eV) & $10^9$, 6.0 & $10^9$, 3.2 & $6.6\times10^{6}$, 2.92\\
particles $\hat\tau_0$ (MPa), $F$ (eV) & -- & -- & 188, 3.66\\
particle climb $D_{c,0}^\dagger$ (m$^2$/s), $Q$ (eV) & -- & -- & $5.4\times10^{-3}$, 2.92\\
\end{tabular}
\end{ruledtabular}
\end{table}

\section{S4. Data sources}
Tungsten: resolved shear stress, $m$ and $V^\star$ from micropillar compression, nanoindentation, indentation literature, IGW, W plate and arc-melted W \todo{references}. Zirconium: resolved shear stress and $V^\star$ from Derep \emph{et al.} \todo{ref}; $m$ from micro-cantilever tests on prismatic and basal slip \todo{Gong ref} and from Lee (1972, 2007) \todo{refs}. Eurofer-97 tensile data: Vanaja \emph{et al.}, J. Nucl. Mater. \textbf{424} (2012) \todo{pages}. Eurofer-97 creep: minimum creep rates digitized from Fig.~1 of Fern\'andez \emph{et al.}~\cite{S_Fern} (plate and bar, air and vacuum, 450--650~$^\circ$C, 53 points), together with the 550~$^\circ$C, 300~MPa steady-state rate of Yu \emph{et al.}~\cite{S_Yu}. The digitization uncertainty is about $\pm2\%$ in stress and $\pm0.05$ decade in rate. The digitized values are provided in \texttt{eurofer97\_creep\_data\_digitized.csv}.

\begin{thebibliography}{9}
\bibitem{S_BrownHam} L. M. Brown and R. K. Ham, in \emph{Strengthening Methods in Crystals}, edited by A. Kelly and R. B. Nicholson (Elsevier, Amsterdam, 1971), p.~9.
\bibitem{S_Fern} P. Fern\'andez, A. M. Lancha, J. Lape\~na, R. Lindau, M. Rieth, and M. Schirra, Fusion Eng. Des. \textbf{75--79}, 1003 (2005).
\bibitem{S_Yu} G. Yu, N. Nita, and N. Baluc, Fusion Eng. Des. \textbf{75--79}, 1037 (2005).
\end{thebibliography}
\end{document}
